% Options for packages loaded elsewhere
% Options for packages loaded elsewhere
\PassOptionsToPackage{unicode}{hyperref}
\PassOptionsToPackage{hyphens}{url}
\PassOptionsToPackage{dvipsnames,svgnames,x11names}{xcolor}
%
\documentclass[
  nonatbib,
  preprint]{article}
\usepackage{xcolor}
\usepackage{amsmath,amssymb}
\setcounter{secnumdepth}{-\maxdimen} % remove section numbering
\usepackage{iftex}
\ifPDFTeX
  \usepackage[T1]{fontenc}
  \usepackage[utf8]{inputenc}
  \usepackage{textcomp} % provide euro and other symbols
\else % if luatex or xetex
  \usepackage{unicode-math} % this also loads fontspec
  \defaultfontfeatures{Scale=MatchLowercase}
  \defaultfontfeatures[\rmfamily]{Ligatures=TeX,Scale=1}
\fi
\usepackage{lmodern}
\ifPDFTeX\else
  % xetex/luatex font selection
\fi
% Use upquote if available, for straight quotes in verbatim environments
\IfFileExists{upquote.sty}{\usepackage{upquote}}{}
\IfFileExists{microtype.sty}{% use microtype if available
  \usepackage[]{microtype}
  \UseMicrotypeSet[protrusion]{basicmath} % disable protrusion for tt fonts
}{}
\makeatletter
\@ifundefined{KOMAClassName}{% if non-KOMA class
  \IfFileExists{parskip.sty}{%
    \usepackage{parskip}
  }{% else
    \setlength{\parindent}{0pt}
    \setlength{\parskip}{6pt plus 2pt minus 1pt}}
}{% if KOMA class
  \KOMAoptions{parskip=half}}
\makeatother
% Make \paragraph and \subparagraph free-standing
\makeatletter
\ifx\paragraph\undefined\else
  \let\oldparagraph\paragraph
  \renewcommand{\paragraph}{
    \@ifstar
      \xxxParagraphStar
      \xxxParagraphNoStar
  }
  \newcommand{\xxxParagraphStar}[1]{\oldparagraph*{#1}\mbox{}}
  \newcommand{\xxxParagraphNoStar}[1]{\oldparagraph{#1}\mbox{}}
\fi
\ifx\subparagraph\undefined\else
  \let\oldsubparagraph\subparagraph
  \renewcommand{\subparagraph}{
    \@ifstar
      \xxxSubParagraphStar
      \xxxSubParagraphNoStar
  }
  \newcommand{\xxxSubParagraphStar}[1]{\oldsubparagraph*{#1}\mbox{}}
  \newcommand{\xxxSubParagraphNoStar}[1]{\oldsubparagraph{#1}\mbox{}}
\fi
\makeatother


\usepackage{longtable,booktabs,array}
\newcounter{none} % for unnumbered tables
\usepackage{calc} % for calculating minipage widths
% Correct order of tables after \paragraph or \subparagraph
\usepackage{etoolbox}
\makeatletter
\patchcmd\longtable{\par}{\if@noskipsec\mbox{}\fi\par}{}{}
\makeatother
% Allow footnotes in longtable head/foot
\IfFileExists{footnotehyper.sty}{\usepackage{footnotehyper}}{\usepackage{footnote}}
\makesavenoteenv{longtable}
\usepackage{graphicx}
\makeatletter
\newsavebox\pandoc@box
\newcommand*\pandocbounded[1]{% scales image to fit in text height/width
  \sbox\pandoc@box{#1}%
  \Gscale@div\@tempa{\textheight}{\dimexpr\ht\pandoc@box+\dp\pandoc@box\relax}%
  \Gscale@div\@tempb{\linewidth}{\wd\pandoc@box}%
  \ifdim\@tempb\p@<\@tempa\p@\let\@tempa\@tempb\fi% select the smaller of both
  \ifdim\@tempa\p@<\p@\scalebox{\@tempa}{\usebox\pandoc@box}%
  \else\usebox{\pandoc@box}%
  \fi%
}
% Set default figure placement to htbp
\def\fps@figure{htbp}
\makeatother


% definitions for citeproc citations
\NewDocumentCommand\citeproctext{}{}
\NewDocumentCommand\citeproc{mm}{%
  \begingroup\def\citeproctext{#2}\cite{#1}\endgroup}
\makeatletter
 % allow citations to break across lines
 \let\@cite@ofmt\@firstofone
 % avoid brackets around text for \cite:
 \def\@biblabel#1{}
 \def\@cite#1#2{{#1\if@tempswa , #2\fi}}
\makeatother
\newlength{\cslhangindent}
\setlength{\cslhangindent}{1.5em}
\newlength{\csllabelwidth}
\setlength{\csllabelwidth}{3em}
\newenvironment{CSLReferences}[2] % #1 hanging-indent, #2 entry-spacing
 {\begin{list}{}{%
  \setlength{\itemindent}{0pt}
  \setlength{\leftmargin}{0pt}
  \setlength{\parsep}{0pt}
  % turn on hanging indent if param 1 is 1
  \ifodd #1
   \setlength{\leftmargin}{\cslhangindent}
   \setlength{\itemindent}{-1\cslhangindent}
  \fi
  % set entry spacing
  \setlength{\itemsep}{#2\baselineskip}}}
 {\end{list}}
\usepackage{calc}
\newcommand{\CSLBlock}[1]{\hfill\break\parbox[t]{\linewidth}{\strut\ignorespaces#1\strut}}
\newcommand{\CSLLeftMargin}[1]{\parbox[t]{\csllabelwidth}{\strut#1\strut}}
\newcommand{\CSLRightInline}[1]{\parbox[t]{\linewidth - \csllabelwidth}{\strut#1\strut}}
\newcommand{\CSLIndent}[1]{\hspace{\cslhangindent}#1}



\setlength{\emergencystretch}{3em} % prevent overfull lines

\providecommand{\tightlist}{%
  \setlength{\itemsep}{0pt}\setlength{\parskip}{0pt}}



 


\makeatletter
\@ifpackageloaded{caption}{}{\usepackage{caption}}
\AtBeginDocument{%
\ifdefined\contentsname
  \renewcommand*\contentsname{Table of contents}
\else
  \newcommand\contentsname{Table of contents}
\fi
\ifdefined\listfigurename
  \renewcommand*\listfigurename{List of Figures}
\else
  \newcommand\listfigurename{List of Figures}
\fi
\ifdefined\listtablename
  \renewcommand*\listtablename{List of Tables}
\else
  \newcommand\listtablename{List of Tables}
\fi
\ifdefined\figurename
  \renewcommand*\figurename{Figure}
\else
  \newcommand\figurename{Figure}
\fi
\ifdefined\tablename
  \renewcommand*\tablename{Table}
\else
  \newcommand\tablename{Table}
\fi
}
\@ifpackageloaded{float}{}{\usepackage{float}}
\floatstyle{ruled}
\@ifundefined{c@chapter}{\newfloat{codelisting}{h}{lop}}{\newfloat{codelisting}{h}{lop}[chapter]}
\floatname{codelisting}{Listing}
\newcommand*\listoflistings{\listof{codelisting}{List of Listings}}
\makeatother
\makeatletter
\makeatother
\makeatletter
\@ifpackageloaded{caption}{}{\usepackage{caption}}
\@ifpackageloaded{subcaption}{}{\usepackage{subcaption}}
\makeatother
\usepackage{bookmark}
\IfFileExists{xurl.sty}{\usepackage{xurl}}{} % add URL line breaks if available
\urlstyle{same}
\hypersetup{
  pdftitle={The Attention Cycle: Reducing Supervision Tax through Executable Documents},
  pdfauthor={Lianghui Zhang},
  pdfkeywords={position paper, supervision tax, document-mediated
collaboration, intent-bound work, human-AI collaboration, attention
management, Getting Things Done},
  colorlinks=true,
  linkcolor={blue},
  filecolor={Maroon},
  citecolor={Blue},
  urlcolor={Blue},
  pdfcreator={LaTeX via pandoc}}


\title{The Attention Cycle: Reducing Supervision Tax through Executable
Documents}
\author{Lianghui Zhang}
\date{}
\begin{document}
\maketitle
\begin{abstract}
\emph{This is a position paper.} Current agent systems pursue autonomy
by making agents run longer with less human input. For well-defined,
verifiable tasks this is the right direction. But for
\textbf{intent-bound work} --- work whose value depends on deciding
\emph{what} should be done, not only on doing it correctly --- the
bottleneck is not agent autonomy but intent specification. When intent
is under-specified, autonomy does not remove supervision; it converts
the human into a runtime specifier who clarifies, redirects, approves,
and repairs agent work as ambiguities surface during execution. We call
this continuous attention cost the \textbf{supervision tax} of
under-specified intent, and it scales with agent parallelism while human
attention does not.

This paper proposes \textbf{document-mediated collaboration} as a way to
cut that tax without removing the human from the loop. The human and a
foreground AI use Deep Focus sessions to think the work through; for the
executable slice, they externalize intent, constraints, context, and
definitions of done into durable task documents. Once clarified, these
documents become executable contracts for background agents, which work
asynchronously and write results back. The human re-enters through a
Decision Queue and periodic Review rather than continuous supervision.

We call this operating rhythm the \textbf{Attention Cycle}: Deep Focus →
Parallel Execution → Decision Queue, supported by an independent Review
perspective that prevents the cycle from collapsing into an epistemic
echo chamber. The cycle treats human attention as the scarce, fragile
resource it is --- invested in thinking, protected during execution, and
re-engaged only at decision boundaries. Empirical validation of
supervision-tax reduction is left to future work.
\end{abstract}


\subsection{1. Introduction}\label{introduction}

\subsubsection{1.1 The wrong problem: removing the
human}\label{the-wrong-problem-removing-the-human}

Current agent systems are largely organized around one ambition: make
agents more autonomous. They should plan further, run longer, recover
from failures, coordinate with tools, and require less human input. For
routine, verifiable, well-specified tasks, this is the right direction.
If the task is clear, the success condition is objective, and the cost
of error is bounded, the best interface is often no interface at all.

But many tasks are not like this. Research, product design,
architecture, writing, strategy, and complex engineering belong to what
we will call \textbf{intent-bound work}: work whose value depends on
\emph{what should be done}, not only on whether the doing is technically
correct. In intent-bound work, deciding what to attempt, what to leave
out, and what ``good'' means is itself most of the value. The human is
not in the loop because they type commands; they are in the loop because
they supply intent, judgment, taste, context, and accountability.
Removing them does not save labor --- it removes the source of value.

The real problem is therefore not how to remove the human. It is how to
remove the wrong kind of human involvement.

The human should not be supervising every intermediate step, answering
clarification questions at runtime, approving partial outputs, and
repairing agent drift. But the human must remain deeply engaged where
their judgment matters: specifying intent, choosing direction,
evaluating tradeoffs, and deciding whether the result is good enough.

This paper argues that the key bottleneck in serious human--AI work is
not agent autonomy, but \textbf{supervision tax}: the continuous
attention cost imposed on the human when agents run on intent that was
not specified well enough to run without them.

The goal is not to move the human out of the loop. The goal is to move
the human out of runtime supervision.

\subsubsection{1.2 Supervision tax comes from under-specified
intent}\label{supervision-tax-comes-from-under-specified-intent}

A common explanation for supervision is that agents are still not
capable enough. This is partly true, but incomplete. Even very capable
agents require supervision when the task they receive is
under-specified.

For an agent to act autonomously, it must make decisions along the way.
If the task does not carry enough intent, context, constraints, and
definition of done, the agent fills the gaps by inference. Some
inferences are harmless. Others encode assumptions the human never
agreed to. As execution proceeds, these hidden assumptions surface as
clarification requests, partial misalignments, unexpected choices, and
results that are technically competent but directionally wrong.

At that point, the human becomes a runtime specifier. They clarify what
should have been clarified earlier, redirect the agent after it has
already deviated, approve fragments whose purpose is still unstable, and
repair outputs generated from implicit assumptions. The agent may be
autonomous in the narrow sense that it can keep taking actions. But the
work is not autonomous, because the human must keep re-entering the
process to supply missing specification.

This gives the first half of the paper's opening claim:

\begin{quote}
Autonomy does not eliminate supervision when intent is under-specified;
it relocates it from the agent's runtime into the human's attention.
\end{quote}

The second half follows from how autonomy scales. One agent on
under-specified intent produces a stream of clarification requests,
partial outputs, and decision points flowing back to the human. Ten
agents in parallel produce ten such streams converging on the same
attention. Parallelism does not absorb under-specification; it
concentrates it into more concurrent surface area than one human can
keep up with:

\begin{quote}
The more parallel the agents, the more under-specified intent amplifies
into supervision tax --- clarification load scales with agents; human
attention does not.
\end{quote}

Supervision tax is therefore not a UI problem nor a temporary model
limitation. It is the runtime cost of skipped clarification, scaled by
parallelism.

\subsubsection{1.3 AI-assisted clarification and executable
documents}\label{ai-assisted-clarification-and-executable-documents}

The frontier of agent design is moving in exactly the direction that the
previous section made worrying: toward longer-running, more autonomous
agents. The ambition is to run further on each invocation, recover from
more failures, coordinate more tools, and call back to the human less
often. For routine work, this is genuine progress. For intent-bound
work, however, it raises the bar on what the agent is given to start
with. The longer and more autonomously an agent operates, the more
weight its initial specification has to carry --- every gap, ambiguity,
or unstated constraint becomes either an unsupervised wrong turn or a
runtime callback to the human. \textbf{More-autonomous agents do not
weaken the demand for clarified intent; they sharpen it.}

Producing clarified intent of that quality is itself difficult, and it
is where chat-form AI is at its strongest. A capable foreground AI is a
\emph{thinking partner} --- exposing ambiguity, pushing back on vague
claims, forcing precision --- which makes the intent-bound work that
resists agent automation paradoxically well-served by AI as a discussion
medium. The pattern this paper proposes is to use that strength
deliberately: the human and a foreground AI think the work through
together, and materialize the parts ready to commit as durable documents
that background agents can act on. §2 develops the document side; §3
develops the attention rhythm around it.

This gives the operating pattern this paper proposes:

\begin{quote}
\textbf{Document-mediated collaboration.} A human and a foreground AI
clarify intent together in conversation; the result of that
clarification is materialized as an executable document; background
agents act on the document asynchronously and write results back; the
human re-enters through review and decision rather than continuous
supervision.
\end{quote}

Chat is where thinking happens. Documents are where clarified intent
runs.

\subsubsection{1.4 The Attention Cycle}\label{the-attention-cycle}

§1.3 located the irreplaceable human contribution at \emph{upstream
clarification}. That contribution has an awkward property: it does not
scale. Clarification cannot be parallelized across many humans without
losing the coherence that gives it value; it cannot be sped up by
pushing harder. It is bound to a fixed, fragile resource --- the human's
attention. Document-mediated collaboration removes runtime supervision;
it does not remove the \emph{upstream} attention cost of clarification
itself. The remaining lever is to make that attention as effective as
possible.

Attention is scarce in two senses, and the cycle has to address both.

It is scarce in \textbf{quantity}. Recall §1.2's closing line:
\emph{clarification load scales with agents; human attention does not.}
As more agents run in parallel, more clarification work converges on one
fixed attention budget. Throwing capacity at the problem does not work,
because there is only one head.

It is scarce in \textbf{quality}, in a way pure budgeting misses. The
deep, sustained focus that clarification work depends on takes effort to
enter, is easily disrupted, and is slow to rebuild. The real cost of an
interruption is not the minute spent answering it; it is the focus state
that has to be reconstructed afterward --- and the flow state, if any,
that is lost in the process. The payoff of \emph{protected} attention is
therefore non-linear: one uninterrupted hour routinely outperforms two
half-hours fragmented by clarification pings, and a sustained flow state
can outperform either by a wide margin.

Making upstream clarification efficient is therefore an
attention-management problem. The cycle has two jobs: spend the human's
limited attention only where it adds value, and \emph{protect} the
uninterrupted focus state under which that attention is actually worth
spending.

The inspiration comes from \textbf{Getting Things Done}: externalize
open loops into a trusted system so nothing competes for attention
except what genuinely needs it now. Document-mediated collaboration
extends that move into AI-mediated work --- once intent is materialized
into an executable document, the running task lives in the document and
the agents acting on it, not in the human's head.

The \textbf{Attention Cycle} turns these principles into an operating
rhythm:

\begin{verbatim}
Deep Focus → Parallel Execution → Decision Queue → Deep Focus
\end{verbatim}

with \textbf{Review} as a structurally separate attention slot that
challenges the cycle from outside. Each phase is defined by the
\emph{kind} of attention it requires --- or deliberately does not:
\textbf{Deep Focus} is undivided attention spent thinking with a
foreground AI; \textbf{Parallel Execution} requires none;
\textbf{Decision Queue} batches judgment on the human's terms;
\textbf{Review} is a context-independent perspective that keeps the
cycle from collapsing into an echo chamber. §3 develops each phase in
turn.

\subsubsection{1.5 Contributions}\label{contributions}

Against the diagnosis developed above --- that serious human--AI work
pays a \textbf{supervision tax} whenever autonomous agents run on
under-specified intent --- this paper makes two constructive
contributions.

First, it proposes \textbf{document-mediated collaboration} as the
mechanism for cutting that tax: chat with a foreground AI handles the
thinking, executable documents handle the running work, and the human
re-enters through review and decision rather than continuous
supervision. §2 develops it.

Second, it defines the \textbf{Attention Cycle} as the operating rhythm
that organizes the human's attention across \textbf{Deep Focus},
\textbf{Parallel Execution}, \textbf{Decision Queue}, and an independent
\textbf{Review} perspective, so that agent capability can keep scaling
without scaling demands on the human. §3 develops it.

Taken together, these contributions are an attempt at what we take to be
the \textbf{actual productivity bottleneck of the current AI era}: the
capable-agent side of human--AI work has advanced extraordinarily fast;
the efficient-human side has not. This paper is an attempt to close that
gap from the human side.

\subsection{\texorpdfstring{\textbf{2. From Supervision Tax to
Executable
Documents}}{2. From Supervision Tax to Executable Documents}}\label{from-supervision-tax-to-executable-documents}

\subsubsection{\texorpdfstring{\textbf{2.1 Supervision tax as the cost
of weak
specification}}{2.1 Supervision tax as the cost of weak specification}}\label{supervision-tax-as-the-cost-of-weak-specification}

For intent-bound work, stronger agents do not reduce the need for
specification. They increase the consequence of getting specification
wrong.

A routine task can often be executed from a short instruction because
much of the intent is implicit in the task type. If the goal is to
resize an image, run a test suite, summarize a document, or file a
ticket, the agent can rely on established conventions and objective
completion criteria. But intent-bound work is different. The difficult
part is not only executing the task, but deciding what the task should
mean: what direction matters, what constraints are binding, what
tradeoffs are acceptable, and what ``good'' means.

When that intent is not made explicit, the agent fills the gaps by
inference. A weak agent may fail early. A stronger agent can proceed
further, making more decisions along the way, each one based on
assumptions the human may never have authorized. The result can be
competent execution in the wrong direction.

This is the supervision tax of under-specified intent. The human must
re-enter during execution to clarify, redirect, approve, or repair what
should have been specified before the agent started.

The important point is not that agents are incapable. It is that
autonomy raises the standard for the input artifact. The more capable
and long-running the agent becomes, the more the initial task
description has to carry.

\begin{verbatim}
more autonomous agents → higher specification burden
under-specified intent → runtime supervision tax
\end{verbatim}

This is why reducing supervision tax requires better upstream
clarification, not merely stronger downstream execution.

\subsubsection{\texorpdfstring{\textbf{2.2 Thinking with a foreground
AI}}{2.2 Thinking with a foreground AI}}\label{thinking-with-a-foreground-ai}

If stronger agents raise the specification burden, the next question is
how high-quality specifications are produced. The answer turns out to be
indirect. A good specification is not the thing the human is producing
during a Deep Focus session --- it is a \emph{byproduct} of something
larger. What is actually being produced is a deeper understanding of the
work itself, of which the specification is only the externalized
residue.

For intent-bound work, the human rarely begins with a complete view of
the problem already formed in their head. They begin with a direction, a
concern, a partial judgment about what might matter, a hunch they have
not yet examined. Most of the work is not transcribing intent that
already exists; it is \emph{forming} it. To do that, the human has to
surface implicit constraints, test framings, reconsider assumptions,
separate what is decided from what is still open, and figure out what
``good'' even means in this case. This is thinking --- slow, structured,
often uncomfortable --- and it is what intent-bound work actually
consists of.

This is also where chat-form AI is at its strongest. A capable
foreground AI is a \emph{thinking partner}: it can hold context across a
long line of reasoning, propose alternative framings, push back on vague
claims, raise considerations the human has not yet weighed, and force
precision where it is being avoided. Its primary value is not that it
writes good specifications. Its primary value is that it makes the
\emph{human's own thinking} sharper. The human leaves a good Deep Focus
session understanding the problem better than they did going in --- and
that deepened understanding is the real product. Specifications are one
downstream form of it, not the point of it.

The interaction is therefore:

\begin{verbatim}
human and AI think together
 → human’s understanding deepens
 → clarified intent emerges as one product among others
\end{verbatim}

Some of what emerges from these sessions is internal and never
externalized at all: a shift in how the human sees the problem, a
revised intuition, a discarded direction, a quietly upgraded mental
model. Some is externalized but not executable on its own: macro plans,
design notes, framings, decisions worth recording, syntheses of
reference material. And some --- when a slice of the thinking has
converged enough to commit --- becomes a specification precise enough
that another agent can act on it. The framework cares centrally about
this last category, because it is what enables asynchronous execution;
but it is important to be honest that it is a \emph{subset} of what Deep
Focus produces, not the whole of it.

\subsubsection{\texorpdfstring{\textbf{2.3 Materializing clarified
intent into executable
documents}}{2.3 Materializing clarified intent into executable documents}}\label{materializing-clarified-intent-into-executable-documents}

A chat session is a good medium for thinking, but a poor substrate for
asynchronous work. It is exploratory, linear, full of discarded
alternatives, and bound to the conversation that produced it. The
durable products of a Deep Focus session have to be concluded and
externalized into the workspace if anything other than the human --- and
in many cases, even the human at a later time --- is going to act on
them.

Not all of those products are task documents. The same document
workspace naturally accumulates a broader \textbf{knowledge substrate}:
notes, plans, framings, decisions, reference syntheses, and other
durable outputs of Deep Focus that are not themselves executable. Task
documents live alongside this substrate, link into it for context, and
are read by background agents against it. The substrate is the
surrounding ground that makes any single task document coherent. The
framework assumes such a substrate exists and is reasonably maintained,
but it does not propose a knowledge-management methodology of its own.
Curation of the substrate is a real concern (§5.1); it is not the
subject of this paper.

Within that substrate, a \textbf{task document} is the \emph{executable
subset} --- the kind of document that turns a clarified slice of intent
into a stable artifact a background agent can act on. It records the
task's intent, context, constraints, accepted decisions, open questions,
definition of done, and later execution results. It becomes the single
source of truth for that piece of work: the background agent reads from
it, the human reviews through it, the Review perspective critiques it,
and execution history accumulates back into it.

This is what streamlines the collaboration. For each piece of executable
work, clarification, execution, and review no longer fragment across
chat threads, agent logs, and side notes that must be reconciled later.
They flow through one artifact:

\begin{verbatim}
conversation → task document → agent execution → writeback → review
\end{verbatim}

The document becomes an \textbf{executable contract} when it is rich
enough for a background agent to act on and explicit enough for
completion to be judged without returning to the human for live
clarification. It does not need to encode everything as rigid fields.
Much of the specification can remain prose. What matters is that the
document --- together with the surrounding knowledge substrate it links
into --- carries enough context, constraints, and acceptance criteria
for the agent to know whether it should complete the task, report a
blocker, or raise a decision request.

This gives the central transformation of the framework:

\begin{verbatim}
clarified intent → durable task document → executable contract
\end{verbatim}

In this sense, document-mediated collaboration is not simply ``using
documents with agents.'' It is a coordination pattern with a sharp
focus: chat is where thinking happens, the knowledge substrate is where
its durable outputs accumulate, task documents are the runnable subset
of that substrate, and Review is where human judgment re-enters without
continuous monitoring.

\subsection{\texorpdfstring{\textbf{3. The Attention
Cycle}}{3. The Attention Cycle}}\label{the-attention-cycle-1}

Section 2 described how Deep Focus thinking produces durable artifacts
--- including, for the executable slice, task documents agents can act
on. This moves specification out of runtime supervision and into
upstream clarification. But it does not remove the upstream cost of the
thinking itself. That thinking still requires human attention, and that
attention remains scarce, fragile, and non-scalable.

The \textbf{Attention Cycle} answers that question:

\begin{verbatim}
Deep Focus → Parallel Execution → Decision Queue → Deep Focus
\end{verbatim}

with \textbf{Review} as a structurally separate attention slot that
provides independent challenge to the cycle.

The cycle is built around a simple principle: anything that does not
require the human's attention right now should not be in it. This
principle echoes Getting Things Done: open loops occupy attention until
they are externalized into a trusted system. In document-mediated
collaboration, executable documents extend that principle to agentic
work. Once intent has been clarified and materialized into a document,
the running task no longer has to live in the human's attention. It can
live in the document and in the agents acting on it.

The cycle therefore separates four attention modes: deep thinking,
zero-attention execution, batched judgment, and independent review.

\subsubsection{\texorpdfstring{\textbf{3.1 Deep Focus: thinking as
investment}}{3.1 Deep Focus: thinking as investment}}\label{deep-focus-thinking-as-investment}

Deep Focus is the phase where the human's attention is deliberately
spent.

The human engages one objective with full attention: defining intent,
choosing direction, evaluating tradeoffs, naming constraints, and
deciding what ``good'' means. A foreground AI participates as a
\emph{thinking partner}. It helps surface hidden assumptions, compare
alternatives, structure context, ask what is missing, push back on vague
claims, and sharpen the human's own understanding of the problem.

The output of Deep Focus is, first, progress in the human's
understanding of the objective --- a sharper view of the problem that
may not be externalized at all. Where the thinking has converged enough
to commit, it also produces durable artifacts in the workspace: notes,
plans, decisions, and, for the executable slice, one or more task
documents that carry enough clarified intent for background agents to
act without live supervision.

What sets Deep Focus apart from ordinary attentive work is
\emph{uninterruption}. The quality of thinking does not scale linearly
with minutes spent on a problem; it scales with whether those minutes
are continuous enough to enter the cognitive states where deep work
actually happens. Two effects on the negative side make this asymmetry
sharp:

\begin{itemize}
\tightlist
\item
  \textbf{Recovery cost.} Empirical studies of knowledge workers find
  that after an interruption, the average time to return to the original
  task at comparable depth is on the order of 23 minutes (Mark et al.,
  2008). The minute spent answering the interruption is a small fraction
  of its real cost.
\item
  \textbf{Attention residue.} Even when the interrupting task is brief,
  part of attention remains attached to the prior context (Leroy, 2009);
  the worker's effective cognitive bandwidth on the resumed task is
  degraded for some time afterward.
\end{itemize}

Flow research adds the positive-side counterpart: when uninterrupted
focus is sustained long enough, the worker enters a state in which
performance on cognitively demanding work is substantially elevated
above baseline (Csikszentmihalyi, 1990). Flow is fragile --- a single
interruption can dissolve it --- and not reliably re-enterable on the
same day.

The compound effect is that Deep Focus is not merely attention-expensive
but attention-fragile, and the value lost to interruption is non-linear
in interruption frequency. Thinking a problem through to coherence ---
and, where ready, externalizing it as an executable contract whose
intent, constraints, context, and acceptance criteria fit together ---
depends on assembling exactly this kind of attention, which is far
easier to break than to rebuild.

The purpose of the rest of the cycle is therefore partly protective.
Parallel Execution, Decision Queue, and Review are designed so that Deep
Focus is not continuously invaded by agent activity.

\subsubsection{\texorpdfstring{\textbf{3.2 Parallel Execution: no
attention
required}}{3.2 Parallel Execution: no attention required}}\label{parallel-execution-no-attention-required}

Parallel Execution is the phase where background agents act on clarified
documents without consuming human attention.

A background agent picks up a task document, reads the specification,
follows its context links, executes the task, verifies what it can, and
writes results back into the same document system. Depending on the
task, the writeback may include changed artifacts, logs, test results,
summaries, blockers, follow-up tasks, or decision requests.

The important design target is not simply that execution is
asynchronous. It is that execution does not create a live supervision
burden.

The human should not have to watch the agent work, answer incidental
questions as they appear, or keep a mental model of every running task.
If the document is sufficient, the agent proceeds. If it is
insufficient, the agent writes back the missing condition as a blocker
or decision request. In either case, the human's attention is not pulled
into the runtime.

This is where agent parallelism becomes useful rather than costly. Many
agents may run at once, but their activity is mediated through
executable documents. Their parallelism does not become parallel
interruption.

The intended relation is:

\begin{verbatim}
agent capacity scales in parallel
human attention remains bounded
\end{verbatim}

\subsubsection{\texorpdfstring{\textbf{3.3 Decision Queue: judgment
without
interruption}}{3.3 Decision Queue: judgment without interruption}}\label{decision-queue-judgment-without-interruption}

No specification is perfect. Some ambiguity will survive clarification.
Some tasks will finish and require human judgment. Some agents will
encounter blockers that genuinely require a decision. These items should
return to the human, but not as interruptions.

The Decision Queue is the channel through which background execution
re-enters human attention.

A decision item should carry enough structure to be handled without
reconstructing the whole task from memory: source task, question,
available options, recommended option if any, reasoning, urgency,
reversibility, and expected impact. The point is not to hide uncertainty
from the human. The point is to present uncertainty at the right
boundary.

Items enter the queue asynchronously. The human engages them
deliberately.

Some items are light: approvals, triage, small corrections, or obvious
accept/reject decisions. These can be handled during low-intensity
intervals. Other items are heavy: architectural tradeoffs, product
direction, research interpretation, or disagreement with an agent's
proposed path. These should be pulled into a Deep Focus session, because
they are not interruptions; they are intent-bound work in their own
right.

This gives the queue two roles.

First, it protects Deep Focus by preventing agents from interrupting
whenever they encounter ambiguity.

Second, it preserves human judgment by ensuring that unresolved
ambiguity still returns to the human, but in a batched and reviewable
form.

The queue is therefore not a supervision dashboard. A dashboard invites
the human to monitor execution. The Decision Queue exists to avoid
monitoring. It is a buffer between agent asynchrony and human attention.

\subsubsection{\texorpdfstring{\textbf{3.4 Independent Review
Perspective: friction without
supervision}}{3.4 Independent Review Perspective: friction without supervision}}\label{independent-review-perspective-friction-without-supervision}

A cycle of human and foreground AI working closely inside the same
context has a specific failure mode. The system can become more coherent
while becoming narrower.

The foreground AI helps the human elaborate the current direction. It
recalls supporting context, sharpens arguments, resolves local
ambiguity, and makes the work feel more internally consistent. This is
useful during Deep Focus. But the same dynamic can create an epistemic
echo chamber. The system may optimize local reasoning while drifting
away from external reality, ignoring untested assumptions, or
reinforcing a mistaken frame.

The Attention Cycle therefore needs an \textbf{independent Review
perspective}: a structurally separate attention slot whose purpose is to
challenge the products of the cycle from \emph{outside} the foreground
working context. Its job is to introduce friction, not to continue the
work --- to ask the questions the cycle is structurally unlikely to ask
itself: what assumption has not been tested, where the work has drifted
from its stated criteria, which failures are repeating, what external
evidence would challenge the current direction, what has become locally
coherent but globally questionable.

What makes Review work is not a particular implementation, but a single
invariant: \textbf{independence from the foreground context that
produced the work}. A reviewer that inherits the foreground AI's
conversation, memory, or framing tends to become a more polite version
of the same voice. As long as that independence is preserved, the
concrete mechanism is deliberately left open. Review may be a single
skeptical agent or several; it may run on a cadence, on workspace
events, or on demand; it may be a human playing the role. As model and
agent capabilities improve, the instruction may shrink to little more
than \emph{``look at this with a critical eye.''} The framework's claim
is about the \emph{position} of Review in the cycle, not its
implementation.

Two constraints follow from that position. First, Review must not
reintroduce runtime supervision: its findings enter the same Decision
Queue as any other decision item, to be engaged on the human's terms
rather than as interrupts. Second, Review must remain a
\emph{perspective} rather than an authority --- its outputs are inputs
to human judgment, not verdicts. With those constraints in place, Review
can introduce cognitive friction at the right boundary without breaking
the rest of the cycle.

\subsubsection{\texorpdfstring{\textbf{3.5 The cycle as attention
discipline}}{3.5 The cycle as attention discipline}}\label{the-cycle-as-attention-discipline}

The Attention Cycle is not primarily a task workflow. It is an attention
discipline.

It keeps apart two kinds of work that current agent interfaces often
collapse:

\begin{verbatim}
Human attention:
think, decide, judge, review.

Agent execution:
search, implement, verify, summarize, write back.
\end{verbatim}

The human remains deeply in the loop, but not continuously in the
runtime. Their attention is invested upstream in thinking, protected
during execution, reintroduced through batched judgment, and
periodically challenged through independent review.

This is the core post-supervision pattern. The system does not aim to
eliminate human involvement. It aims to place human involvement where it
has the highest leverage.

For intent-bound work, the bottleneck is not how many actions agents can
take. It is how well the human's limited attention can specify, direct,
and judge those actions. The Attention Cycle lets agent capability keep
scaling while keeping human attention bounded.

\subsection{\texorpdfstring{\textbf{4.
Instantiation}}{4. Instantiation}}\label{instantiation}

The Attention Cycle does not require building a new agent platform. Its
two roles --- a foreground AI used as a thinking partner during Deep
Focus, and a background agent that executes clarified task documents
asynchronously --- are both well covered by existing products. A minimal
instantiation needs only three pieces: a \textbf{chat-agent prompt} that
encodes the foreground role (think with the human, materialize converged
slices as task documents), a \textbf{background-agent prompt} that
encodes the execution role (find Ready tasks, execute against the
document, write results back), and a thin \textbf{document protocol}
specifying the small amount of structure the two agents must coordinate
on --- in practice, a \texttt{Status} field (e.g., Ready / In Progress /
Needs Review / Done / Blocked) and an explicit \texttt{Dependencies}
field. Everything else --- intent, context, constraints, definition of
done, writeback notes, decision requests --- stays as ordinary document
prose, written, read, and edited in plain language.

With those three pieces, the cycle can be assembled today from
off-the-shelf components: any document workspace that supports
structured text and version-aware editing, any conversational agent
capable of editing that workspace, and any execution-oriented agent
capable of running asynchronously against the documents and writing
results back --- whether a coding agent, a browser agent, a shell agent,
or some combination. A dedicated product designed end-to-end around the
cycle is also a reasonable path; it is not, however, a prerequisite for
putting the methodology into practice.

In fact, a fragmented version of this pattern is already in widespread
use. Many practitioners today think through problems with one AI tool,
manually copy the result into another tool to execute it, then copy the
execution result back to continue the discussion. The intuition behind
that workflow is exactly the one this paper is articulating: chat is the
right medium for thinking, agents are the right medium for executing,
and the human moves between them. What is missing is the shared
substrate that lets the two sides exchange state directly, the protocol
that lets them coordinate without the human as a courier, and the
operating rhythm that keeps the human's attention from being shredded by
the handoff.

The Attention Cycle is the attempt to streamline that emergent pattern:
replace copy-paste between tools with one document the agents share,
replace ad-hoc supervision with a Decision Queue, replace the assumption
that the foreground agent is also the executor with a clean separation
of roles around the document boundary. The contribution is not any
specific instantiation; it is that this rhythm is already half-emergent
in how people work with AI today, and that pulling it into a single
coherent operating model is what frees the scarce resource --- human
attention --- to do the work only it can do.

\subsection{5. Limitations and
Conclusion}\label{limitations-and-conclusion}

\subsubsection{\texorpdfstring{\textbf{5.1
Limitations}}{5.1 Limitations}}\label{limitations}

The Attention Cycle depends on a strong premise: human attention can be
protected by moving execution out of the runtime loop and by
materializing clarified intent into documents. This premise has several
boundary conditions.

\textbf{First, not all work is ready to become executable.} The
framework applies best when a slice of work has converged enough that
intent, constraints, context, and definition of done can be faithfully
externalized. In early exploration, tacit-knowledge-heavy craft, or
situations where the right question is still forming, premature
documentization can become premature commitment. In those cases,
conversation should remain exploratory, and the output of Deep Focus may
be only a changed human understanding, not an executable task.

\textbf{Second, the cycle protects attention; it does not multiply it.}
Deep Focus remains expensive, fragile, and non-scalable. The framework
can reduce runtime supervision and improve the throughput of clarified
intent, but it cannot remove the human cost of forming intent in the
first place. If the volume of intent-bound work exceeds the human's
capacity for sustained clarification and judgment, agent parallelism
cannot close that gap.

\textbf{Third, executable documents can still fail as specifications.} A
document may be actionable without being correct. It may omit
assumptions, encode weak acceptance criteria, or allow a capable agent
to proceed confidently in the wrong direction. The Decision Queue and
independent Review perspective are recovery mechanisms for this failure,
not proof that the failure has been eliminated.

\textbf{Fourth, the framework depends on epistemic and documentary
discipline.} Review only works if it remains structurally independent
from the foreground context it is meant to challenge. Documents only
work as executable contracts if they remain current, linked, and
trusted. Prompt drift, accumulated shared context, stale decisions, dead
links, and unintegrated writebacks all erode the cycle's foundation. The
Attention Cycle therefore does not remove the need for workspace hygiene
and review design; it makes those disciplines central to the system.

\subsubsection{\texorpdfstring{\textbf{5.2
Conclusion}}{5.2 Conclusion}}\label{conclusion}

The central argument of this paper is that the next bottleneck in
agentic work is not only agent capability, but human attention. As
agents become more autonomous and more parallel, under-specified intent
does not disappear; it returns as supervision tax. The Attention Cycle
reframes this problem: the goal is not to remove the human from the
loop, but to move human attention to the places where it has the highest
leverage.

This reframing suggests a different direction for agent systems. Instead
of treating interaction as a sequence of prompts and responses, systems
should treat clarified intent as a durable artifact. Documents are not
merely records of work; they can become the coordination substrate
through which humans think, agents execute, reviewers challenge, and
decisions accumulate. In this view, the document is the boundary between
human attention and agent parallelism.

The broader implication is that productivity in the AI era may depend
less on how much work agents can do, and more on how well systems
protect the scarce attention that gives that work direction. More
capable agents make this problem sharper, not weaker. Without better
intent formation, autonomy scales execution faster than judgment. With
executable documents and a disciplined attention cycle, agent capacity
can scale while human attention remains bounded.

Future work should therefore study the Attention Cycle not only as a
conceptual model, but as an operational pattern: what document
structures make intent executable, what review mechanisms preserve
independence, what decision queues minimize interruption without hiding
risk, and how much supervision tax can actually be reduced in real
workflows. The practical question is not whether humans stay in the
loop. They do. The question is where in the loop their attention
belongs.

\protect\phantomsection\label{refs}
\begin{CSLReferences}{1}{1}
\bibitem[\citeproctext]{ref-csikszentmihalyi1990flow}
Csikszentmihalyi, Mihaly. 1990. \emph{Flow: The Psychology of Optimal
Experience}. Harper; Row.

\bibitem[\citeproctext]{ref-leroy2009attention}
Leroy, Sophie. 2009. {``Why Is It so Hard to Do My Work? The Challenge
of Attention Residue When Switching Between Work Tasks.''}
\emph{Organizational Behavior and Human Decision Processes} 109 (2):
168--81.

\bibitem[\citeproctext]{ref-mark2008cost}
Mark, Gloria, Daniel Gudith, and Ulrich Klocke. 2008. {``The Cost of
Interrupted Work: More Speed and Stress.''} \emph{Proceedings of the ACM
Conference on Human Factors in Computing Systems (CHI)}, 107--10.

\end{CSLReferences}




\end{document}
